\section{Execution infrastructure and delivery pipeline}
\label{sec:infrastructure}

\subsection{Laboratory topology}

\subsection{Why a laboratory at all}
\label{sec:proxmox}

An audit cannot be measured on a real network. The result would depend on whatever
happens to be exposed that day, could not be repeated, and would not be
comparable between two runs. LANCE therefore audits networks that it has itself
built, from a description, and throws away afterwards.

A scenario is that description: which devices exist, which services they run,
which weaknesses are installed, and which connections are allowed. Deploying it
means materialising those devices as virtual machines on a hypervisor, so that the
audit sees a real network rather than a simulation, and destroying them afterwards
so that the next trial starts from the same state.

\textbf{Proxmox} provides that hypervisor. It is the host on which the scenario
devices are created as LXC containers --- lightweight, unprivileged Linux
containers, each acting as one device of the scenario: a web service, an IoT
service, a gateway, or any other role the scenario describes. The scenario never
runs on the workstation and never touches a real network: it exists only on the
Proxmox host, for the duration of a trial.

\textbf{Ansible} plays the other half. It configures each container from the
scenario description --- installs the service, sets the credentials, injects the
weakness --- and resets it to a clean state before the next trial. Deployment,
configuration and reset are the three operations that make a trial repeatable; the
hypervisor alone would give none of them.

The cycle is therefore: deploy the scenario as containers, configure and inject,
run the audit, destroy. A run that skipped the reset would inherit the previous
trial's state, and its result would no longer be attributable.

\begin{figure}[H]
\centering
\begin{tikzpicture}[
  box/.style={draw, rounded corners=2pt, align=center, font=\footnotesize,
              minimum height=10mm, text width=30mm, inner sep=3pt},
  arr/.style={-{Latex[length=1.6mm]}, font=\scriptsize, draw=black!70},
  node distance=7mm and 12mm,
]
\node[box] (scen) {Scenario\\[-1pt]\scriptsize devices, services,\\[-1pt]\scriptsize weaknesses, links};
\node[box, below=10mm of scen] (ans) {Ansible\\[-1pt]\scriptsize configure, inject, reset};
\node[box, below=10mm of ans] (px) {Proxmox host\\[-1pt]\scriptsize LXC containers:\\[-1pt]\scriptsize web, IoT, gateway};
\node[box, right=22mm of px] (agent) {Audit agent\\[-1pt]\scriptsize plans and runs\\[-1pt]\scriptsize the tests};
\node[box, above=10mm of agent] (eval) {Evaluator\\[-1pt]\scriptsize recomputes from\\[-1pt]\scriptsize the journal};
% flows
\draw[arr] (scen) -- node[left]{\scriptsize deploy} (ans);
\draw[arr] (ans) -- (px);
\draw[arr] (px) -- node[above]{\scriptsize tests, observations} (agent);
\draw[arr] (agent) -- node[right]{\scriptsize journal} (eval);
\draw[arr, dashed] (agent.west) to[bend left=35] node[below]{\scriptsize destroy} (px.east);
\end{tikzpicture}
\caption{The laboratory cycle. A scenario is deployed as virtual devices on the
Proxmox host, configured by Ansible, audited, then destroyed. The evaluator reads
the journal, never the devices.}
\label{fig:labo}
\end{figure}

The execution of LANCE relies on a single Proxmox host and one LXC container that
carries the entire platform. The playbook
\code{benchmarks/ansible/playbooks/deploy\_master.yml} creates, on the Proxmox host
group, an unprivileged LXC container with identifier 200 named \code{nato-master}
(Debian 13, 20~GB, 4096~MB of memory, 2~cores, \code{nesting=1}, \code{onboot 1}).
Two interfaces are attached: \code{eth0} on \code{vmbr0} in DHCP for management,
and \code{eth1} on \code{vmbr1} in \code{192.168.100.200/24} for the scenario test
network. The playbook adds the TUN device authorisation and the mount
\code{/dev/net/tun}, required by Tailscale, then rebuilds an \code{OUTPUT} chain
of \code{iptables}: on \code{eth0} only the local interface, the
\code{ESTABLISHED,RELATED} states, DNS, HTTP, HTTPS, SSH, the Proxmox API (port
8006), UDP 41641 and the apt cache are allowed; everything else is rejected. On
\code{eth1}, outbound traffic is fully allowed. \code{eth1} has no gateway, which
confines the scans to the isolated network.

The Proxmox host name is given as \code{pve-nato} by \code{AGENTS.md},
\code{objectifs/ETAT-COURANT.md} and \code{docs/guides/umons-llm.md}, but neither
the playbook nor \code{docs/guides/development-instances.md} names it: the
versioned Ansible inventory is a template whose address literally equals
\code{<PROXMOX\_IP>}, and the only constant readable in the repository is the
container identifier, 200. The actual configuration lives outside the repository.

\subsection{The three instances}

The script \code{scripts/deployment/deploy-instance.sh} fixes the slot, port,
directory and service, and rejects any other name. The stable service is produced
by the playbook; the development units are written by the script.

\begin{table}[H]
\centering
\caption{Instances deployed on the master container.}
\label{tab:instances}
\begin{tabular}{@{}llll@{}}
\toprule
Slot & Branch & Port & Checkout and service \\
\midrule
\code{main} & \code{main} & 8501 & \code{/opt/nato-smartcity-iot} / \code{nato-fastapi} \\
\code{dev-1} & \code{dev/1} & 8502 & \code{/opt/lance-dev-1} / \code{lance-dev-1} \\
\code{dev-2} & \code{dev/2} & 8503 & \code{/opt/lance-dev-2} / \code{lance-dev-2} \\
\bottomrule
\end{tabular}
\end{table}

Each slot has its own Python environment, its database, its \code{.env} file and
its artefacts. The shared configuration is written into
\code{/var/lib/lance/<slot>.env} and imported by a systemd override; these entries
feed \code{GET /api/environment}, which returns the slot, the branch, the commit
and the activation of the lock. After start-up the script checks that the HTTP
response announces exactly the validated commit and the expected slot, and that
the lock is active.

The common lock is a single file, \code{/var/lib/lance/lab.lock}, opened by
\code{exec 9>>} and then locked by \code{flock -x 9} before any mutation; the
comment is explicit: never delete this file. The same value is passed to the APIs,
which allows the three instances to coordinate. Waiting is cooperative, without
arrival ordering, and kept in memory. A deployment of a development slot checks
that the stable instance announces the lock: an older stable instance would not
participate in it.

\subsection{The runner and the delivery pipeline}

The GitHub Actions runner is installed in \code{/opt/actions-runner}, registered
with the label \code{nato-master} and run as a systemd service.

The pipeline reads as four links. A \code{push} to \code{main}, \code{dev/1} or
\code{dev/2}, or a merge request into these branches, triggers
\code{.github/workflows/benchmark-integrity.yml}. The first job runs
\code{.github/scripts/ci\_scope.py}, which classifies the modified files and
publishes the \code{full} flag. The second builds the image of the isolated worker
and checks the absence of ground truth. The deployment jobs then call
\code{.github/workflows/update-master.yml}, which runs the deployment script with
the slot, the validated SHA, the branch and the remote repository.

\begin{table}[H]
\centering
\caption{Triggers and paths through the pipeline.}
\label{tab:circuit}
\small
\begin{tabular}{@{}p{3.1cm}p{5.6cm}p{4.8cm}@{}}
\toprule
Event & Path & Deployment \\
\midrule
documentation \code{push} & classification to \code{full=false}: diff of the previous commit and whitespace check & none \\
code \code{push} & full tests, scenario and Ansible checks, sandbox, worker image & \code{main} to 8501, \code{dev/*} to 8502 or 8503 \\
Merge request & same path according to classification & none \\
Manual trigger & \code{full=true} forced & chosen slot, or checks only \\
Version tag & \code{full=true} forced & \code{main} only \\
\bottomrule
\end{tabular}
\end{table}

The light path covers the Markdown, LaTeX, PDF and image files under
\code{docs/}, \code{objectifs/} and \code{research/}, as well as a closed list of
presentation files. The classification keeps deletions and the two paths of a
rename; any undetermined event falls back to the full path, and a failure blocks
the deployment. A documentation push therefore restarts no instance.

On the machine side, the update script fetches the SHA and applies a fast-forward
merge for \code{main}, or a detachment after checking for a clean tree for the
development branches; it then verifies that \code{HEAD} equals the validated SHA.
Before that, the service is stopped, the lock acquired, and the instance queried
in order to refuse mutating active or pending work. In the event of a failure after
stopping, the service remains stopped: no automatic rollback to an earlier version
is performed.

\subsection{Corrections to the deployment script}

The history of the script shows a tight sequence of commits on 29 October 2026.
Until then, the script checked the host identity by a hard-coded address and
required the root user; the first commit of that day removes both assumptions.

\begin{table}[H]
\centering
\caption{Successive corrections and cause addressed.}
\label{tab:fixes}
\small
\begin{tabular}{@{}p{1.6cm}p{3.4cm}p{7.4cm}@{}}
\toprule
Commit & Correction & Cause addressed \\
\midrule
\code{6b0a9ff} & address in an environment variable, \code{sudo} re-execution, checks of the stable instance and of the configuration made optional, Codex CLI under \code{/usr/local/bin} & a rebuilt container broke every deployment: identity, privileges and neighbouring instances assumed inherited \\
\code{433e2a3} & identity deduced from the \code{tailscale0} interface, failing a variable & removes the last constant dependent on a rebuild \\
\code{a568c68} to \code{dcd5903} & absolute path, execution via \code{bash}, propagation of the runner name and of the whole deployment environment & \code{sudo} does not resolve a relative path nor preserve the environment \\
\code{150eb8a} & Codex CLI download failure tolerated & under \code{set -e}, the failure stopped the deployment and left the instance stopped \\
\bottomrule
\end{tabular}
\end{table}

Each correction corresponds to a concrete failure mode, not to an aesthetic
refinement. The same identity logic now applies to the benchmark controllers, which
read the Tailscale interface through the shared helper and are covered by a
dedicated test; a laboratory lock variable lets an operator suspend that check
explicitly rather than by disabling the code.

\subsection{The UMONS model and its relay}

The \code{ollama-umons} provider is reached through a local SSH transfer on the
Proxmox host, towards the university cluster (\code{docs/guides/umons-llm.md}). The
path was identified and restored on 4~October 2026, with a successful minimal
generation; the authentication observed is interactive. The restoration command is
documented, as are the listening checks. The guide is explicit: neither start-up
at boot, nor reconnection after the process ends; the administration of the
downstream GPU server remains to be confirmed with the administrator. The
fragility is established in the facts: the relay went down again on 22 October,
and it was the provider pre-flight that named the failure without executing any
arm.

\subsection{The outage of 29 October}

The outages actually traced in the repository for that day are those of the
container rebuild, which rendered obsolete the fixed address and the assumption of
a privileged runner. By contrast, \emph{no mention of a multipath default route on
the laboratory router exists in the repository}: the search on the corresponding
terms returns nothing. This fact, observed in session but not recorded, must not be
presented as documented. It constitutes a gap in the project's memory, since the
same outage will recur at the next router restart.

\subsection{Known fragilities}

\begin{itemize}
\item \textbf{Unreplicated failure points}: single Proxmox host, single container,
single runner. The runner registration tolerating failure, an absent runner does
not block the playbook but blocks every deployment.
\item \textbf{UMONS relay}: no supervised service, interactive authentication,
address still fixed in the examples of the guide although commit
\code{e54da24} claimed to have stopped fixing addresses in the documentation.
Restarting LANCE does not repair the relay.
\item \textbf{Cooperative lock}: without arrival ordering, in memory, not covered
for an Ansible script launched directly, an old checkout or a manual
intervention. This is not security isolation between collaborators.
\item \textbf{Manual recovery}: a failure after stopping leaves the service
stopped, with no automatic rollback.
\item \textbf{Host identity}: without a declared variable, the identity depends on
the presence of any address on the Tailscale interface; the script does not verify
that this host is indeed the one that should host the instance.
\item \textbf{Partial documentation}: the name of the Proxmox host and the
physical segment correspond only imperfectly to the playbook; the versioned
inventory is a template.
\end{itemize}
